\chapter{C Compilation Overview}

\section{Introduction}

When you write a C program, you're creating human-readable source code. But computers don't execute C code directly—they execute machine instructions, binary patterns that the CPU understands. The \textbf{compilation pipeline} is the bridge between these two worlds, transforming your source code through a series of well-defined stages, each producing intermediate artifacts.

Understanding this pipeline is essential for:
\begin{itemize}
    \item \textbf{Debugging build failures}: Knowing which stage failed helps pinpoint the problem
    \item \textbf{Optimizing performance}: Optimization happens at specific stages
    \item \textbf{Understanding errors}: Different stages produce different kinds of errors
    \item \textbf{Cross-compilation}: Each stage may run on different machines
    \item \textbf{Build system design}: Makefiles and other build tools orchestrate these stages
\end{itemize}

This chapter provides a complete map of the compilation pipeline, showing how C source code becomes an executable program and what tools make this transformation possible.

\section{The Complete Compilation Pipeline}

The C compilation pipeline consists of four main stages, each converting one file format into another:

\begin{verbatim}
┌──────────────────────────────────────────────────────────────────────────┐
│                    Complete Compilation Pipeline                         │
├──────────────────────────────────────────────────────────────────────────┤
│                                                                          │
│   Source Code   Preprocessing   Preprocessed   Compilation   Assembly    │
│   (.c files) ────────> (.i files) ────────> (.s files) ────────>         │
│                                                                          │
│      │                        │                    │                     │
│      │                        │                    │                     │
│      ▼                        ▼                    ▼                     │
│   ┌──────┐                ┌──────┐              ┌──────┐                 │
│   │ .c   │  ──gcc -E──>   │ .i   │  ──gcc -S──> │ .s   │                 │
│   └──────┘                └──────┘              └──────┘                 │
│                                                                          │
│   Assembly        Assembly        Object Files     Linking     Executable│
│   (.s files) ───────────> (.o files) ───────────────────>                │
│                                                                          │
│      │                                 │                    │            │
│      │                                 │                    │            │
│      ▼                                 ▼                    ▼            │
│   ┌──────┐                        ┌──────┐              ┌──────┐         │
│   │ .s   │  ──as -o──>           │ .o   │  ──ld──>      │(none)│         │
│   └──────┘                        └──────┘              └──────┘         │
│                                                         (a.out)          │
└──────────────────────────────────────────────────────────────────────────┘
\end{verbatim}

Each stage has a specific purpose and produces a distinct artifact:

\begin{center}
\begin{tabular}{|l|l|l|l|l|}
\hline
\textbf{Stage} & \textbf{Input} & \textbf{Output} & \textbf{Tool} & \textbf{Purpose} \\
\hline
Preprocessing & \texttt{.c} & \texttt{.i} & \texttt{cpp} (via \texttt{gcc -E}) & Macro expansion, file inclusion \\
\hline
Compilation & \texttt{.i} & \texttt{.s} & \texttt{cc1} (via \texttt{gcc -S}) & Parse, optimize, generate assembly \\
\hline
Assembly & \texttt{.s} & \texttt{.o} & \texttt{as} & Encode to machine code \\
\hline
Linking & \texttt{.o} & executable & \texttt{ld} & Combine objects, resolve symbols \\
\hline
\end{tabular}
\end{center}

\begin{quote}
\textbf{Key Insight}: While we usually invoke \texttt{gcc} or \texttt{clang} to go directly from \texttt{.c} to executable, internally these drivers invoke separate programs for each stage. We can stop at any intermediate stage to inspect the artifacts.
\end{quote}

\subsection{Why Separate Stages Matter}

The modular design of the compilation pipeline serves several important purposes:

\begin{enumerate}
    \item \textbf{Separation of Concerns}: Each stage focuses on one transformation, making the toolchain easier to develop and maintain

    \item \textbf{Optimization Opportunities}: Different optimization techniques apply at different stages. For example:
    \begin{itemize}
        \item Preprocessor macros are expanded before compilation
        \item Code optimization happens during compilation
        \item Link-time optimization (LTO) can happen across translation units
    \end{itemize}

    \item \textbf{Debugging}: When something goes wrong, you can inspect intermediate files to understand where the problem occurred

    \item \textbf{Language Interoperability}: Different frontends can produce the same intermediate representations, allowing code in different languages to be linked together

    \item \textbf{Cross-compilation}: Each stage can theoretically run on different machines (e.g., compiling on a powerful build server for a target embedded system)
\end{enumerate}

Let's now examine each stage in detail.

\section{Stage 1: Preprocessing (\texttt{.c} $\to$ \texttt{.i})}

\textbf{Preprocessing} is the first stage of compilation. It transforms your source code through \textbf{textual substitution} before the compiler ever sees it. The preprocessor doesn't understand C syntax—it works purely on text manipulation.

\subsection{What the Preprocessor Does}

The preprocessor performs three main operations:

\begin{enumerate}
    \item \textbf{File Inclusion}: Inserting the contents of other files (via \texttt{\#include})
    \item \textbf{Macro Expansion}: Replacing macros with their definitions (via \texttt{\#define})
    \item \textbf{Conditional Compilation}: Including or excluding code based on conditions (via \texttt{\#if}, \texttt{\#ifdef}, etc.)
\end{enumerate}

Consider this example:

\begin{lstlisting}[language=C]
// main.c
#define MAX_SIZE 100
#define SQUARE(x) ((x) * (x))

#include <stdio.h>

#ifdef DEBUG
    #define LOG(msg) printf("DEBUG: \%s\n", msg)
#else
    #define LOG(msg) /* do nothing */
#endif

int main(void) {
    int arr[MAX_SIZE];
    LOG("Program started");
    int result = SQUARE(5);
    return 0;
}
\end{lstlisting}

After preprocessing (run \texttt{gcc -E main.c -o main.i}), we get a massive file (typically thousands of lines) where:
\begin{itemize}
    \item \texttt{\#include <stdio.h>} has been replaced with the entire contents of \texttt{stdio.h} (and all headers it includes)
    \item \texttt{MAX\_SIZE} has been replaced with \texttt{100}
    \item \texttt{SQUARE(5)} has been replaced with \texttt{((5) * (5))}
    \item The \texttt{LOG} macro has been replaced either with \texttt{printf} or nothing, depending on whether \texttt{DEBUG} is defined
\end{itemize}

\subsection{Preprocessing Example}

Let's create a simpler example to see the transformation clearly:

\begin{lstlisting}[language=C]
// example.c
#define PI 3.14159
#define AREA(r) (PI * (r) * (r))

double circle_area(double radius) {
    return AREA(radius);
}
\end{lstlisting}

Running \texttt{gcc -E example.c} produces:

\begin{lstlisting}[language=C]
# 1 "example.c"
# 1 "<built-in>"
# 1 "<command-line>"
# 1 "example.c"

double circle_area(double radius) {
    return (3.14159 * (radius) * (radius));
}
\end{lstlisting}

Notice that:
\begin{itemize}
    \item All \texttt{\#define} directives are gone
    \item The macros have been expanded
    \item Preprocessor directives (lines starting with \texttt{\#}) show the original source location for debugging
\end{itemize}

\subsection{Translation Units}

The output of preprocessing is called a \textbf{translation unit}—a single, complete source file ready for compilation. Each \texttt{.c} file becomes one translation unit, which the compiler processes independently.

This is crucial for understanding compilation:
\begin{itemize}
    \item Each \texttt{.c} file is compiled \textbf{separately}
    \item The compiler doesn't see other \texttt{.c} files
    \item Cross-file references are resolved during \textbf{linking}
\end{itemize}

\begin{quote}
\textbf{Common Pitfall}: Because each \texttt{.c} file is compiled independently, you must ensure definitions are consistent across translation units. This is why we use header files (\texttt{.h} files)—to share declarations between compilation units.
\end{quote}

We'll explore preprocessing in depth in Chapter 2.

\section{Stage 2: Compilation (\texttt{.i} $\to$ \texttt{.s})}

\textbf{Compilation} proper transforms the preprocessed source code into assembly language. This is where the heavy lifting happens: parsing, semantic analysis, optimization, and code generation.

\subsection{What the Compiler Does}

The compiler performs several sophisticated operations:

\begin{enumerate}
    \item \textbf{Lexical Analysis}: Breaking the code into tokens (identifiers, keywords, operators)
    \item \textbf{Parsing}: Building an Abstract Syntax Tree (AST) based on C grammar
    \item \textbf{Semantic Analysis}: Type checking, ensuring variables are declared before use, etc.
    \item \textbf{Intermediate Representation}: Converting the AST into an IR for optimization
    \item \textbf{Optimization}: Applying various transformations to improve performance
    \item \textbf{Code Generation}: Emitting assembly language for the target architecture
\end{enumerate}

\subsection{Compilation Example}

Let's compile a simple function:

\begin{lstlisting}[language=C]
// add.c
int add(int a, int b) {
    return a + b;
}
\end{lstlisting}

First preprocess it:
\begin{verbatim}
gcc -E add.c -o add.i
\end{verbatim}

Then compile to assembly:
\begin{verbatim}
gcc -S add.i -o add.s -fno-asynchronous-unwind-tables -fno-stack-protector
\end{verbatim}

The resulting assembly (on x86-64 Linux):

\begin{lstlisting}[language=nasm]
    .intel_syntax noprefix
    .globl  add
    .type   add, @function
add:
    mov eax, edi
    add eax, esi
    ret
\end{lstlisting}

Even from this simple example, we can see:
\begin{itemize}
    \item The function is marked as global (\texttt{.globl})
    \item Arguments are in registers (\texttt{edi} and \texttt{esi} per the calling convention)
    \item The return value goes in \texttt{eax}
    \item The \texttt{ret} instruction returns to the caller
\end{itemize}

\subsection{Why Assembly?}

You might wonder: why generate assembly instead of going directly to machine code? Several reasons:

\begin{enumerate}
    \item \textbf{Readability}: Assembly is human-readable, useful for debugging
    \item \textbf{Portability}: The same compiler can use different assemblers for different targets
    \item \textbf{Flexibility}: Developers can write hand-optimized assembly when needed
    \item \textbf{Debugging}: Assembly listings with source annotations help understand what the compiler did
\end{enumerate}

\subsection{Optimization Levels}

The compiler's behavior changes dramatically based on optimization levels:

\begin{verbatim}
gcc -O0 add.c -o add.o    # No optimization (default for debugging)
gcc -O1 add.c -o add.o    # Basic optimization
gcc -O2 add.c -o add.o    # Standard optimization (recommended)
gcc -O3 add.c -o add.o    # Aggressive optimization
gcc -Os add.c -o add.o    # Optimize for code size
\end{verbatim}

With \texttt{-O2}, our \texttt{add} function becomes even simpler (the compiler might inline it entirely if it's small enough).

We'll explore compilation internals in depth in Chapter 3.

\section{Stage 3: Assembly (\texttt{.s} $\to$ \texttt{.o})}

\textbf{Assembly} converts human-readable assembly language into machine code, producing an \textbf{object file}. This is a binary format, but not yet executable.

\subsection{What the Assembler Does}

The assembler:

\begin{enumerate}
    \item \textbf{Parses Assembly}: Converts mnemonic instructions into their binary encodings
    \item \textbf{Resolves Symbols}: Creates symbol table entries for labels and external references
    \item \textbf{Generates Sections}: Places code and data into appropriate sections (\texttt{.text}, \texttt{.data}, etc.)
    \item \textbf{Creates Relocations}: Records locations that need fixing during linking
\end{enumerate}

\subsection{Assembly Example}

Let's assemble our \texttt{add.s} file:

\begin{verbatim}
as add.s -o add.o
\end{verbatim}

The resulting \texttt{add.o} is a binary file in \textbf{ELF format} (Executable and Linkable Format). We can inspect it:

\begin{verbatim}
$ file add.o
add.o: ELF 64-bit LSB relocatable, x86-64, version 1 (SYSV), not stripped

$ ls -l add.o
-rw-r--r-- 1 user user 864 Nov 15 10:30 add.o
\end{verbatim}

The file is only 864 bytes, containing:
\begin{itemize}
    \item Machine code for the \texttt{add} function
    \item A symbol table noting that \texttt{add} is a global symbol
    \item Section headers marking where code and data are located
    \item Relocation information (none needed for this simple function)
\end{itemize}

\subsection{Object File Contents}

Object files contain several key components:

\begin{center}
\begin{tabular}{|l|l|}
\hline
\textbf{Component} & \textbf{Purpose} \\
\hline
\textbf{Section Headers} & Describe the sections (\texttt{.text}, \texttt{.data}, etc.) \\
\hline
\textbf{Code Sections} & Machine code (\texttt{.text} for executable code) \\
\hline
\textbf{Data Sections} & Initialized data (\texttt{.data}), zero-initialized data (\texttt{.bss}) \\
\hline
\textbf{Symbol Table} & List of defined and undefined symbols \\
\hline
\textbf{Relocation Entries} & Places needing address fixups during linking \\
\hline
\textbf{Debug Information} & (optional) Debug symbols for gdb \\
\hline
\end{tabular}
\end{center}

\subsection{Why Object Files?}

Object files are \textbf{relocatable}—they don't have fixed memory addresses. This allows:
\begin{itemize}
    \item Multiple object files to be linked together
    \item Code to be loaded at any memory address
    \item Libraries to be shared across multiple programs
\end{itemize}

The linker will ultimately resolve all symbols and assign final addresses.

We'll explore assembly and object files in depth in Chapter 4.

\section{Stage 4: Linking (\texttt{.o} $\to$ executable)}

\textbf{Linking} is the final stage, combining multiple object files and libraries into a single executable program. This is where separate compilation units finally come together.

\subsection{What the Linker Does}

The linker performs several critical tasks:

\begin{enumerate}
    \item \textbf{Symbol Resolution}: Connecting references to definitions across object files
    \item \textbf{Relocation}: Assigning final memory addresses and fixing references
    \item \textbf{Library Linking}: Including code from standard and user libraries
    \item \textbf{Executable Creation}: Setting up the executable file format with proper headers
\end{enumerate}

\subsection{Linking Example}

Let's create a complete program with two files:

\begin{lstlisting}[language=C]
// main.c
extern int add(int a, int b);

int main(void) {
    int result = add(5, 3);
    return result == 8 ? 0 : 1;
}
\end{lstlisting}

\begin{lstlisting}[language=C]
// add.c
int add(int a, int b) {
    return a + b;
}
\end{lstlisting}

Compile and assemble separately:

\begin{verbatim}
gcc -c main.c -o main.o    # Stops after assembly
gcc -c add.c -o add.o      # Stops after assembly
\end{verbatim}

Now link them:

\begin{verbatim}
gcc main.o add.o -o program
\end{verbatim}

Or using the linker directly:

\begin{verbatim}
ld main.o add.o -o program \
   /usr/lib/x86_64-linux-gnu/crt1.o \
   /usr/lib/x86_64-linux-gnu/crti.o \
   /usr/lib/x86_64-linux-gnu/crtn.o \
   -lc -dynamic-linker /lib64/ld-linux-x86-64.so.2
\end{verbatim}

The \texttt{gcc} version is simpler because it automatically includes:
\begin{itemize}
    \item \textbf{C runtime startup code} (\texttt{crt1.o}, \texttt{crti.o}, \texttt{crtn.o})
    \item \textbf{Standard library} (\texttt{-lc} for libc)
    \item \textbf{Dynamic linker} path
\end{itemize}

\subsection{Symbol Resolution}

Symbol resolution is the linker's most critical job. Consider our example:

\begin{verbatim}
main.o:  UNDEFINED: add
         DEFINED:   main

add.o:   DEFINED:   add
\end{verbatim}

The linker sees that \texttt{main.o} references \texttt{add}, finds it defined in \texttt{add.o}, and connects them. If \texttt{add} were not defined anywhere, the linker would error:

\begin{verbatim}
ld: main.o: in function 'main':
main.c:(.text+0x14): undefined reference to 'add'
\end{verbatim}

\subsection{Static vs. Dynamic Linking}

Linking can happen at different times:

\textbf{Static Linking}:
\begin{itemize}
    \item All library code is included in the executable
    \item Executables are larger but self-contained
    \item No runtime dependencies
    \item Command: \texttt{gcc program.o -o program -static} or \texttt{gcc program.o -o program -Wl,-Bstatic -lfoo}
\end{itemize}

\textbf{Dynamic Linking} (default):
\begin{itemize}
    \item Library code is referenced but not included
    \item Executables are smaller
    \item Libraries loaded at runtime
    \item Allows library updates without recompiling
    \item Command: \texttt{gcc program.o -o program} (default)
\end{itemize}

\begin{verbatim}
$ ldd program  # Shows dynamic library dependencies
    linux-vdso.so.1
    libc.so.6 => /lib/x86_64-linux-gnu/libc.so.6
    /lib64/ld-linux-x86-64.so.2
\end{verbatim}

\subsection{Library Search Order}

The linker searches for libraries in a specific order:

\begin{enumerate}
    \item Paths specified with \texttt{-L} flags (left to right)
    \item Standard system library directories (\texttt{/usr/lib}, \texttt{/lib}, etc.)
    \item Paths in \texttt{LD\_LIBRARY\_PATH} environment variable (at runtime)
\end{enumerate}

\textbf{Important}: The order of \texttt{-l} flags matters! Libraries are searched left-to-right, and a library is only used to resolve undefined symbols from libraries to its left.

\begin{verbatim}
# Wrong: libfoo needs symbols from libbar, but libbar is searched first
gcc program.o -o program -lfoo -lbar

# Correct: libbar is available when libfoo needs it
gcc program.o -o program -lbar -lfoo
\end{verbatim}

We'll explore linking and relocation in depth in Chapter 6.

\section{Toolchain Components and Roles}

Now that we've seen the complete pipeline, let's identify the actual tools involved:

\subsection{The GCC Toolchain}

When you install \texttt{gcc}, you actually get a collection of programs:

\begin{verbatim}
$ which gcc cc1 as ld
/usr/bin/gcc
/usr/lib/gcc/x86_64-linux-gnu/11/cc1
/usr/bin/as
/usr/bin/ld
\end{verbatim}

\begin{center}
\begin{tabular}{|l|l|l|l|}
\hline
\textbf{Tool} & \textbf{Full Name} & \textbf{Role} & \textbf{Stage} \\
\hline
\texttt{gcc}/\texttt{clang} & Compiler Driver & Orchestrates the pipeline & All \\
\hline
\texttt{cpp} & C Preprocessor & Macro expansion, includes & Preprocessing \\
\hline
\texttt{cc1} & C Compiler & Parses, optimizes, generates assembly & Compilation \\
\hline
\texttt{as} & Assembler & Assembly to machine code & Assembly \\
\hline
\texttt{ld} & Linker & Combines objects, resolves symbols & Linking \\
\hline
\end{tabular}
\end{center}

\begin{quote}
\textbf{Note}: \texttt{gcc} is technically a \textbf{compiler driver}—it doesn't do compilation itself. Instead, it invokes the appropriate tools (\texttt{cpp}, \texttt{cc1}, \texttt{as}, \texttt{ld}) in the correct order with the right arguments.
\end{quote}

\subsection{The Clang Toolchain}

Clang follows a similar architecture:

\begin{center}
\begin{tabular}{|l|l|}
\hline
\textbf{Tool} & \textbf{Role} \\
\hline
\texttt{clang} & Compiler driver (front-end) \\
\hline
\texttt{llvm-as} & LLVM assembler \\
\hline
\texttt{llc} & LLVM static compiler \\
\hline
\texttt{lld} & LLVM linker \\
\hline
\end{tabular}
\end{center}

\subsection{Compiler Drivers: Why They Exist}

You might wonder: if we have separate tools (\texttt{cpp}, \texttt{cc1}, \texttt{as}, \texttt{ld}), why do we typically just run \texttt{gcc}?

\textbf{Compiler drivers} exist to:
\begin{enumerate}
    \item \textbf{Simplify}: Hide complexity from the user
    \item \textbf{Standardize}: Provide a consistent interface across platforms
    \item \textbf{Orchestrate}: Run tools in the correct order
    \item \textbf{Configure}: Add necessary flags and include paths automatically
\end{enumerate}

When you run:

\begin{verbatim}
gcc -O2 -Wall program.c -o program
\end{verbatim}

The driver:
\begin{enumerate}
    \item Determines the file type (C source)
    \item Constructs the command line for \texttt{cpp}
    \item Constructs the command line for \texttt{cc1} (adding \texttt{-O2} optimization)
    \item Constructs the command line for \texttt{as}
    \item Constructs the command line for \texttt{ld} (adding necessary libraries)
    \item Runs each tool sequentially
    \item Cleans up intermediate files (unless you use \texttt{-save-temps})
\end{enumerate}

\subsection{Stopping at Intermediate Stages}

You can stop the compilation process at any stage:

\begin{verbatim}
# Stop after preprocessing (keep .i file)
gcc -E program.c -o program.i

# Stop after compilation (keep .s file)
gcc -S program.c -o program.s

# Stop after assembly (keep .o file)
gcc -c program.c -o program.o

# Keep all intermediate files
gcc -save-temps program.c -o program
# Creates: program.i, program.s, program.o, program
\end{verbatim}

This is incredibly useful for:
\begin{itemize}
    \item \textbf{Debugging}: Seeing what the preprocessor or compiler actually did
    \item \textbf{Learning}: Understanding how C code transforms to assembly
    \item \textbf{Optimizing}: Inspecting compiler output at different optimization levels
\end{itemize}

\subsection{Multiple Source Files}

For projects with multiple source files:

\begin{verbatim}
# Compile all files separately (faster for incremental builds)
gcc -c main.c -o main.o
gcc -c util.c -o util.o
gcc -c data.c -o data.o
gcc main.o util.o data.o -o program

# Or let the driver do it (one command)
gcc main.c util.c data.c -o program
\end{verbatim}

The second approach compiles each \texttt{.c} file to \texttt{.o}, then links them all.

\section{A Complete Example: Tracing the Pipeline}

Let's trace a complete program through all four stages:

\begin{lstlisting}[language=C]
// greet.c
#include <stdio.h>

#define GREETING "Hello, World!"

int main(void) {
    printf("%s\n", GREETING);
    return 0;
}
\end{lstlisting}

\subsection{Stage 1: Preprocessing}

\begin{verbatim}
gcc -E greet.c -o greet.i
\end{verbatim}

\texttt{greet.i} contains the preprocessed source (over 2000 lines due to \texttt{stdio.h} expansion). The key parts:

\begin{lstlisting}[language=C]
# 1 "greet.c"
# 1 "<built-in>"
# 1 "<command-line>"
# 1 "greet.c"
# 1 "/usr/include/stdio.h" 1 3 4
... [thousands of lines from stdio.h] ...
# 2 "greet.c" 2

int main(void) {
    printf("%s\n", "Hello, World!");
    return 0;
}
\end{lstlisting}

Notice:
\begin{itemize}
    \item \texttt{\#include <stdio.h>} is gone, replaced by its contents
    \item \texttt{GREETING} macro has been expanded to \texttt{"Hello, World!"}
\end{itemize}

\subsection{Stage 2: Compilation}

\begin{verbatim}
gcc -S greet.i -o greet.s
\end{verbatim}

\texttt{greet.s} contains assembly code. Simplified version:

\begin{lstlisting}[language={[x86masm]Assembler}]
    .file   "greet.c"
    .intel_syntax noprefix
    .section    .rodata
.LC0:
    .string "Hello, World!"
    .text
    .globl  main
    .type   main, @function
main:
    push    rbp
    mov     rbp, rsp
    sub     rsp, 16
    lea     rdi, [rip + .LC0]
    call    printf
    mov     eax, 0
    leave
    ret
    .size   main, .-main
\end{lstlisting}

Key observations:
\begin{itemize}
    \item The string literal is in the \texttt{.rodata} section (read-only data)
    \item \texttt{printf} is called with the string address
    \item Return value 0 is placed in \texttt{eax}
\end{itemize}

\subsection{Stage 3: Assembly}

\begin{verbatim}
as greet.s -o greet.o
\end{verbatim}

\texttt{greet.o} is now a binary object file:

\begin{verbatim}
$ file greet.o
greet.o: ELF 64-bit LSB relocatable, x86-64, version 1 (SYSV), not stripped

$ size greet.o
   text    data     bss     dec     hex filename
    135       8       0     143      8f greet.o
\end{verbatim}

The object file contains:
\begin{itemize}
    \item 135 bytes of code (\texttt{.text})
    \item 8 bytes of data (the string in \texttt{.rodata})
    \item A symbol table marking \texttt{main} as defined and \texttt{printf} as undefined
    \item Relocation entries for the string reference and \texttt{printf} call
\end{itemize}

\subsection{Stage 4: Linking}

\begin{verbatim}
gcc greet.o -o greet
# Or explicitly:
ld greet.o -o greet \
   /usr/lib/x86_64-linux-gnu/crt1.o \
   /usr/lib/x86_64-linux-gnu/crti.o \
   /usr/lib/x86_64-linux-gnu/crtn.o \
   -lc -dynamic-linker /lib64/ld-linux-x86-64.so.2
\end{verbatim}

Now \texttt{greet} is an executable:

\begin{verbatim}
$ file greet
greet: ELF 64-bit LSB executable, x86-64, version 1 (SYSV), dynamically linked, interpreter /lib64/ld-linux-x86-64.so.2

$ ls -l greet
-rwxr-xr-x 1 user user 16856 Nov 15 10:45 greet

$ ./greet
Hello, World!
\end{verbatim}

\section{Key Takeaways}

\begin{enumerate}
    \item \textbf{The C compilation pipeline has four stages}: Preprocessing, Compilation, Assembly, and Linking. Each stage transforms one file format into another.

    \item \textbf{Preprocessing (\texttt{.c} $\to$ \texttt{.i})}: Textual substitution through macro expansion, file inclusion, and conditional compilation. Creates translation units.

    \item \textbf{Compilation (\texttt{.i} $\to$ \texttt{.s})}: Parses code, builds an AST, generates intermediate representation, applies optimizations, and emits assembly language.

    \item \textbf{Assembly (\texttt{.s} $\to$ \texttt{.o})}: Converts assembly to machine code, creating relocatable object files with symbol tables and relocations.

    \item \textbf{Linking (\texttt{.o} $\to$ executable)}: Combines object files, resolves symbols, applies relocations, and creates the final executable with proper runtime setup.

    \item \textbf{Toolchain components}: \texttt{gcc}/\texttt{clang} are compiler drivers that orchestrate \texttt{cpp}, \texttt{cc1}, \texttt{as}, and \texttt{ld}. Understanding each tool helps debug build issues.

    \item \textbf{Separate compilation}: Each \texttt{.c} file becomes one translation unit, compiled independently. Cross-file references are resolved during linking.

    \item \textbf{You can stop at any stage}: Use \texttt{-E} (preprocessor), \texttt{-S} (compilation), \texttt{-c} (assembly), or \texttt{-save-temps} (keep everything) to inspect intermediate artifacts.

    \item \textbf{Optimization levels matter}: \texttt{-O0} (fast compilation, good for debugging), \texttt{-O2} (recommended for production), \texttt{-O3} (maximum optimization), \texttt{-Os} (optimize for size).

    \item \textbf{Understanding the pipeline helps}: Debug build failures, optimize performance, interpret compiler errors, and design effective build systems.
\end{enumerate}

\section{Looking Ahead}

Now that we've seen the complete compilation pipeline, we'll examine each stage in detail:

\begin{itemize}
    \item \textbf{Chapter 2}: Preprocessing—how macros work, include file handling, conditional compilation
    \item \textbf{Chapter 3}: Compilation internals—lexers, parsers, ASTs, IR, optimization
    \item \textbf{Chapter 4}: Assembly and object files—assembler syntax, object file structure
    \item \textbf{Chapter 5}: ELF format—headers, sections, segments, symbols, relocations
    \item \textbf{Chapter 6}: Linking and relocation—symbol resolution, static vs. dynamic linking
\end{itemize}
